%%
% 引言或背景
% 引言是论文正文的开端,应包括毕业论文选题的背景、目的和意义;对国内外研究现状和相关领域中已有的研究成果的简要评述;介绍本项研究工作研究设想、研究方法或实验设计、理论依据或实验基础;涉及范围和预期结果等。要求言简意赅,注意不要与摘要雷同或成为摘要的注解。
% modifier: 黄俊杰(huangjj27, 349373001dc@gmail.com)
% update date: 2017-04-15
%%

\chapter{绪论}
%定义，过去的研究和现在的研究，意义，与图像分割的不同,going deeper
\label{cha:introduction}
\section{研究背景与意义}
\label{sec:background}
% What is the problem
% why is it interesting and important
% Why is it hards, why do naive approaches fails
% why hasn't it been solved before
% what are the key components of my approach and results, also include any specific limitations，do not repeat the abstract
%contribution
\subsection{研究背景}

自20世纪中期集成电路（Integrated Circuit, IC）诞生以来，半导体技术经历了持续数十年的高速发展，逐步推动人类社会进入信息化与智能化时代。从早期的简单逻辑电路到如今包含数百亿晶体管的复杂系统级芯片（System-on-Chip, SoC），集成电路在计算能力、能效比以及功能集成度等方面均实现了跨越式提升。这一发展过程在很大程度上遵循摩尔定律（Moore's Law），即单位面积芯片上可容纳的晶体管数量随时间呈指数级增长。近年来，随着工艺节点逐步推进至 7nm、5nm 甚至 3nm，芯片设计复杂度已达到前所未有的水平。以现代高性能处理器为例，其晶体管数量已达到数百亿量级，内部连接关系高度复杂，设计难度急剧上升。

然而，随着器件尺寸不断缩小、集成电路规模不断增大，传统依赖人工经验的设计方式已难以应对如此庞大且复杂的系统设计需求。在此背景下，电子设计自动化（Electronic Design Automation, EDA）技术应运而生，并逐渐发展成为集成电路产业链中不可或缺的核心支撑技术。EDA 工具通过将电路设计问题抽象为数学建模与算法优化问题，借助计算机程序自动完成电路设计的各个环节，显著提升了设计效率与设计质量。从硬件描述语言（Hardware Description Language, HDL）建模，到逻辑综合、时序分析，再到物理设计与验证，EDA 技术贯穿芯片设计全流程。

在 EDA 设计流程中，根据设计抽象层次的不同，通常可划分为前端设计（Front-End）与后端设计（Back-End）。前端设计主要关注电路的功能实现与逻辑描述，而后端设计则负责将抽象逻辑结构映射为实际物理版图。随着工艺节点不断推进，后端设计在整个芯片设计流程中的重要性日益凸显，其复杂度在许多场景下已超过前端设计，成为影响芯片性能与可制造性的关键环节。

在后端设计流程中，布局（Placement）与布线（Routing）是计算复杂度很高、工程影响很显著的两个重要步骤。布局阶段的目标是将电路中的标准单元（Standard Cells）合理放置在芯片版图上，使其在满足设计规则前提下尽可能优化线长、密度与时序性能；布线阶段则在布局结果基础上，将所有逻辑单元通过金属互连线连接起来，形成完整电路结构，布线问题由于涉及大量变量与复杂约束，是典型的 NP 难问题。为降低问题复杂度，布线过程通常分为全局布线（Global Routing）与详细布线（Detailed Routing）两个阶段。全局布线在粗粒度层面为每个线网规划大致路径并分配布线资源；详细布线则在精细几何层面生成满足工艺规则的具体走线。布局与布线质量直接决定芯片的性能（Performance）、功耗（Power）以及面积（Area），即 PPA 指标，是衡量芯片设计优劣的核心标准。

% 其中，。在这一过程中，需要综合考虑线长、通孔数量、拥塞程度以及时序约束等多个目标，属于典型的多目标优化问题。

围绕本文关注的两个核心课题，全局布线与布局分别面临不同但相互耦合的挑战：

对于全局布线而言，首先是质量挑战：随着设计规模与互连复杂度持续提升，线网间资源竞争显著加剧，拥塞热点更易形成并在迭代中累积，导致获得高质量全局解的难度不断上升。尤其在先进工艺节点下，线宽、线距、跨层切换与制造约束共同压缩可行解空间，使传统的串行全局布线算法容易陷入局部最优。其次是计算效率挑战：现代设计常包含数百万乃至上亿个标准单元与数千万线网，全局布线中的候选评估、代价更新与冲突消解计算量迅速膨胀，传统基于中央处理器（Central Processing Unit, CPU）的串行流程在运行时间与内存开销方面逐渐逼近工程上限，难以满足工业界对快速迭代的需求。
此外，随着计算机体系结构持续演进，并行计算逐渐成为突破全局布线速度瓶颈的重要技术路径。特别是图形处理单元（Graphics Processing Unit, GPU）在高吞吐数据处理方面具有显著优势，为大规模候选评估、拥塞代价计算与批量路径更新等任务提供了可行加速基础。然而，布线过程仍存在复杂的数据依赖与资源竞争关系，如何在保证优化稳定性和结果正确性的前提下实现高效并行，依然是需要重点攻克的问题。

对于布局环节，核心挑战主要来自极为复杂的约束建模与真实后端结果之间的鸿沟。真实布局场景包含大量复杂且难以统一建模的约束，因此现有布局算法通常需要简化约束、近似代价或连续化处理，以保证可解性与运行效率。这类方法能够较快给出粗粒度可行解，但其优化目标与后端工具真实评价之间仍存在偏差，导致布局结果仍保留一定可挖掘的增量空间。在超大规模设计中，为控制求解复杂度，布局流程可先以模块或分组为优化单元，优先确定模块级边界与相对位置，而不直接进行实例级精细微调。本文将这一阶段性结果称为粗粒度模块级布局规划结果。这种结果既能继承既有模块化方法在收敛速度和全局结构上的优势，又保留了可进一步优化的空间，并且能够作为与主流后端工具低耦合衔接的统一接口。基于这一认识，本文选择在已有模块级布局规划之上开展反馈驱动的增量优化研究，通过细粒度方式对布局约束进行微调，以期获得更高质量的布局优化结果。%，引入多段模拟退火（Simulated Annealing, SA）与受控随机扰动机制，对约束进行小幅细粒度调整，从而在可控计算预算下进一步提升布局质量与后续可布线性。

% 此外，传统布局与布线流程通常采用先布局、后布线的串行设计模式，两者之间缺乏充分的信息交互机制。这种割裂式流程使布局阶段难以准确预估布线需求，从而可能产生严重拥塞问题，而这些问题往往需要在后续布线阶段通过代价高昂的撤销重布（Rip-up and Reroute）操作修复，不仅增加计算开销，也限制最终设计质量的提升。因此，如何在布局阶段引入布线感知机制，实现布局与布线的协同优化，已成为当前 EDA 研究的重要方向。


综上所述，在先进工艺节点与超大规模设计背景下，全局布线面临质量和速度的权衡，布局环节则面临不充分建模与真实物理设计结果之间的鸿沟。本文据此从两条主线展开研究：一是在全局布线侧构建动态拥塞感知与并行加速的初始布线工作流；二是在布局侧针对粗粒度结果引入多段模拟退火细化优化流程，以实现后端设计质量与效率的协同提升。

\subsection{研究意义}

针对上述背景，开展基于并行计算及模块级优化的 VLSI 高效布局布线算法研究具有重要理论价值与工程意义。

从EDA工具求解质量的层面出发，布局布线问题本质上属于大规模组合优化问题，涉及图论、运筹学与数值优化等多个领域：

\begin{enumerate}
	\item 在全局布线阶段，通过引入拥塞感知的并行布线的优化思想，可从新的视角对传统问题进行建模与求解。例如，将复杂的全局布线优化问题建模成统一的代价表达式，从而可以考虑全局的资源利用情况。
	\item 在布局阶段，通过引入细粒度的优化算法，可以在模块级布局规划结果基础上对布局引导约束进行进一步微调，从而一定程度弥补解析式粗粒度布局算法与超大设计空间搜索之间的鸿沟。例如在模块级布局规划结果之上引入多段模拟退火算法，通过设计高效的随机邻域扰动机制，结合既有的后端设计工具真实布局结果提供的反馈信息可高效地对布局结果进一步微调，以获得更好的布局优化结果。
\end{enumerate}
% 在全局布线阶段，通过引入拥塞感知的并行布线的优化思想，可从新的视角对传统问题进行建模与求解。例如，将复杂的全局布线优化问题建模成统一的代价表达式，从而可以考虑全局的资源利用情况。

% 在布局阶段，通过引入细粒度的优化算法，可以在模块级布局规划结果基础上对布局引导约束进行进一步微调，从而一定程度弥补解析式粗粒度布局算法与超大设计空间搜索之间的鸿沟。例如在模块级布局规划结果之上引入多段模拟退火算法，通过设计高效的随机邻域扰动机制，结合既有的后端设计工具真实布局结果提供的反馈信息可高效的对布局结果进一步微调，以获得更好的布局优化结果。

% 其次，在方法层面，将布局与布线进行协同建模具有重要意义。通过在布局阶段引入可布线性（Routability）约束，可在设计早期主动规避潜在拥塞问题，从源头降低布线难度。这种前瞻性优化思想相比传统事后修复策略，更有利于获得高质量设计结果。同时，通过动态拥塞感知机制，可在布线过程中实时调整优化方向，使算法具备更强全局适应能力，从而提升整体优化效果。

从计算层面出发，在计算密集的全局布线任务中，若可以充分利用并行计算平台的算力优势，对于提升全局布线工具性能具有重要意义。通过设计适用于GPU的并行算法，可显著加速布线过程并缩短芯片设计周期。更短设计周期通常意味着更快产品上市速度，具有一定的经济意义。此外，GPU具有的数据高吞吐量的特点有利于支持更大规模电路设计，为未来超大规模集成电路布线算法技术提供一定参考价值。

从工程应用层面出发，在约束极为复杂的布局环节，传统算法为了提升运行的速度，因此在建模的时候往往更改一些约束设定，例如简化一些约束或者连续化离散的约束条件等等以支持GPU的加速任务。除此之外，真实的物理制造场景中也存在一些非常规且难建模的约束，因此传统的布局算法所得到的布局结果仍存在一定的提升空间，且此提升空间难以通过单纯设计更复杂的布局算法来获得。
% 从工程应用角度看，高效布局布线算法能够直接提升芯片设计质量。在拥塞控制方面，通过优化单元布局与布线路径，可有效减少资源冲突并提高布线成功率；在线长优化方面，可降低信号传输延迟并提升芯片性能；在功耗控制方面，通过减少不必要绕线与通孔，可降低动态功耗与静态功耗。这些改进对于高性能计算、人工智能芯片与移动终端等应用场景均具有重要意义。

% 最后，从技术发展趋势看，随着人工智能（Artificial Intelligence, AI）与 EDA 技术融合不断加深，未来布局布线算法将更加依赖高效数据处理与并行计算能力。本文围绕并行计算与模块化优化展开研究，不仅可为当前 EDA 工具性能提升提供支撑，也为未来 AI 驱动的 EDA 方法奠定基础。

综上所述，本文意图针对全局布线和布局的环节展开研究，在不同的环节探索面向超大规模集成电路的高效布局布线算法，因此该研究课题对于提升后端设计工具的求解质量与运算性能具有一定的意义。%该研究对于提升 EDA 工具性能、缩短芯片设计周期以及推动集成电路产业发展具有重要意义。同时，本文在理论建模与方法设计层面的探索，也可为后续相关研究提供有价值的参考。

\section{研究内容}
\label{sec:research_content}

围绕第1.1节所述规模、拥塞与串行瓶颈等关键挑战，本文将研究内容划分为两个环节：全局布线环节与布局环节。前者聚焦于在高复杂度线网场景下提升布线质量的同时权衡计算效率，后者聚焦于弥补粗粒度布局建模与真实后端结果之间的差距，在既有模块级布局规划结果基础上挖掘增量优化空间。

本文主要采用以下两种思路。

第一，全局布线环节的质量-效率协同优化思路。针对传统流程中后期拥塞修复代价高、局部最优显著以及串行求解效率受限的问题，本文在初始布线阶段引入动态拥塞感知机制与候选树更新策略，通过全局代价解析建模增强早期决策质量。在数据级并行与线网级并行的层面重构关键流程，以提升超大规模集成电路全局布线场景下对布线资源的综合权衡能力并利用GPU加速布线流程。

第二，布局环节中的细粒度增量优化思路。针对布局环节中复杂约束难以充分建模、粗粒度解析结果与真实后端评价存在偏差的问题，本文在已有模块级布局规划结果基础上引入多段模拟退火（Simulated Annealing, SA）流程，通过高效随机扰动与工具反馈驱动的迭代优化，对布局约束进行小幅细粒度微调，在可控预算下进一步提升布局质量.



基于上述思路，本文形成的具体研究工作如表\ref{tab:chap1_contributions}所示。

\begin{table}[htbp]
	\centering
	\caption{本文研究工作与预期作用}
	\label{tab:chap1_contributions}
	\begin{tabular}{p{0.14\textwidth}p{0.30\textwidth}p{0.46\textwidth}}
		\toprule
		研究环节 & 研究工作 & 预期作用 \\
		\midrule
		\multirow{2}{*}{全局布线} & 质量提升：动态拥塞感知初始布线工作流 & 抑制早期拥塞扩散，降低后续撤销重布与迷宫布线开销。 \\
		 & 速度提升：面向并行平台的全局布线加速 & 提升超大规模设计的计算效率与可扩展性，使耗时不随设计规模增长而爆炸。 \\
		\midrule
		布局 & 模块级布局规划细粒度微调：利用多段模拟退火实现增量优化 & 在既有模块级布局规划结果基础上进行小幅优化，改善布局质量。 \\
		\bottomrule
	\end{tabular}
\end{table}

% 需要说明的是，本文在布局方向的工作聚焦于基于模块级布局规划结果的外层优化流程与协同机制；模块级布局核心建模与核心优化算法属于既有研究基础，本文主要关注其在完整布局-布线流程中的集成应用与增量收益验证。该边界设定保证了研究贡献表述的准确性，也与后续章节展开保持一致。

\section{国内外研究现状和相关工作}
\label{sec:related_work}

超大规模集成电路物理设计直接决定芯片性能、功耗与面积等核心指标。物理设计主要包含布局与布线两大阶段：布局为模块与标准单元提供物理位置基础，布线在此基础上完成电气互连。随着工艺节点向 3nm 及以下演进、晶体管规模迈入百亿级，传统基于 CPU 的串行布局布线算法在效率、可扩展性与全局优化能力方面面临显著挑战。围绕这一问题，研究范式已从传统串行组合优化逐步演进到并行计算加速、模块级建模等优化方向。

本节将先讨论与本文研究内容重点最直接相关的全局布线研究，再讨论面向布线质量提升的布局相关研究。

为便于横向比较不同研究路线，表\ref{tab:routing_compare}与表\ref{tab:placement_compare}分别给出全局布线与布局方向代表性方法的对比。

\begin{table}[htbp]
	\centering
	\caption{全局布线方向代表性研究路线对比}
	\label{tab:routing_compare}
	\begin{tabular}{p{0.20\textwidth}p{0.21\textwidth}p{0.23\textwidth}p{0.24\textwidth}}
		\toprule
		路线类别 & 代表工作 & 主要优势 & 主要局限 \\
		\midrule
		组合式全局布线 & FastRoute\cite{10.1145/1233501.1233596}、FastRoute 4.0\cite{10.5555/1509633.1509768} & 工程成熟度高，流程完整。 & 早期拥塞感知有限，串行瓶颈明显。 \\
		\midrule
		并行与多线程布线 & NCTU-GR 2.0\cite{Liu2013NCTUGR2M}、SPRoute 2.0\cite{10.1109/ASP-DAC52403.2022.9712557}、GAMER\cite{9643563} & 能有效缩短部分阶段运行时间。 & 阶段间耦合弱，全流程并行协同不足。 \\
		\midrule
		可微与异构布线 & DGR\cite{10.1145/3649329.3656530}、FastGR\cite{10.24963/ijcai.2023/720}、动态拥塞感知初始布线\cite{11017940} & 具备全局联合优化潜力，适配 GPU 并行。 & 对树候选质量、冲突消解与稳定性仍有提升空间。 \\
		\bottomrule
	\end{tabular}
\end{table}

\begin{table}[htbp]
	\centering
	\caption{布局方向代表性研究路线对比}
	\label{tab:placement_compare}
	\begin{tabular}{p{0.20\textwidth}p{0.20\textwidth}p{0.23\textwidth}p{0.25\textwidth}}
		\toprule
		路线类别 & 代表工作 & 主要优势 & 主要局限 \\
		\midrule
		随机优化布局 & TimberWolf\cite{1586125} & 全局搜索能力较强，早期效果稳定。 & 计算复杂度高，难以扩展到超大规模设计。 \\
		\midrule
		解析式布局 & GORDIAN\cite{Kleinhans1991GORDIANVP}、Kraftwerk2\cite{10.1109/TCAD.2008.925783}、EPlace\cite{10.1145/2699873}、RePlAce\cite{10.1109/TCAD.2018.2859220} & 收敛速度快，适配大规模优化。 & 对复杂约束耦合与跨阶段协同支持有限。 \\
		\midrule
		并行与学习驱动布局 & DREAMPlace\cite{9122053}、DREAMPlace 3.0\cite{9256736}、GraphPlanner\cite{10.1145/3555804} & 计算效率高，适合异构平台与大规模数据。 & 统一的代价建模方法缺少一定自定义空间 \\
		\bottomrule
	\end{tabular}
\end{table}

\subsection{VLSI 全局布线算法研究现状与进展}

\subsubsection{全局布线核心流程与基础研究}

全局布线通常由生成树、模板布线与迷宫布线三个阶段构成。生成树阶段常采用 FLUTE\cite{10.1109/TCAD.2007.907068} 快速构建直角斯坦纳最小树（Rectilinear Steiner Minimal Tree, RSMT）；模板布线阶段对两引脚子问题进行模板选择与层分配；迷宫布线阶段针对拥塞或冲突区域进行精细化修复。

FastRoute\cite{10.1145/1233501.1233596} 通过拥塞驱动树构建与高效模式布线，推动了布线前移思想。后续 FastRoute 4.0\cite{10.5555/1509633.1509768} 进一步关注通孔优化。与此同时，面向特殊约束的树生成与避障研究持续推进，例如高效避障 RSMT 方法\cite{Guo2024ARH}、浅光树优化路线\cite{10.5555/3199700.3199776} 以及学习驱动树构造方法\cite{10.1145/3394885.3431566}。

\subsubsection{全局布线各阶段方法演进}

在树生成阶段，传统仅以线长为主的策略容易忽略后续拥塞，限制全局解质量。围绕该问题，研究者提出了拥塞感知树构建、快速拥塞估计及并行生成策略。针对并行化瓶颈，GPU 加速 RSMT 生成工作\cite{10.1145/3508352.3549434} 证明了数据级并行在早期阶段的有效性。

在模板布线阶段，研究路线经历了从二维近似到三维联合优化的演进。NCTU-GR 2.0\cite{Liu2013NCTUGR2M}、SPRoute 2.0\cite{10.1109/ASP-DAC52403.2022.9712557} 等方法在多线程与详细可布线性约束方面进行了系统改进。近年来，基于 DAG 的全局布线框架（如 EDGE\cite{10.1109/DAC56929.2023.10247702}）与可微全局布线（DGR\cite{10.1145/3649329.3656530}）进一步增强了全局视角下的联合优化能力。本文相关工作也在该方向上提出动态拥塞感知的解析初始布线流程\cite{11017940}，以提升前期拥塞抑制能力。

在迷宫布线阶段，算法重点在于搜索效率与解质量平衡。NCTU-GR 2.0\cite{Liu2013NCTUGR2M}、SPRoute 2.0\cite{10.1109/ASP-DAC52403.2022.9712557} 在有界搜索与并行调度方面提供了代表性方案。GPU 路线中，早期 GPU 全局路由器\cite{6657028} 与 GAMER\cite{9643563} 分别从 A* 搜索与并行扫描角度实现加速，验证了该阶段的并行潜力。

\subsubsection{并行计算在全局布线中的应用}

传统全局布线器多以 CPU 串行为主，在超大规模场景下常出现运行时间过长与可扩展性不足问题。为此，近年来研究逐步转向 CPU-GPU 协同与异构并行。FastGR\cite{10.24963/ijcai.2023/720} 通过异构任务图调度实现 CPU-GPU 协同加速，体现了数据级并行 + 任务级并行融合的有效性。

从方法机制看，并行全局布线主要包括两个层面：其一，数据级并行，即对拥塞图更新、代价计算、候选路径评估等高吞吐任务进行并行映射；其二，线网级并行，即在冲突控制与一致性约束下并行处理多个线网。两者协同有望同时提升运行效率与全局解质量。

\subsubsection{全局布线现有研究局限性}

尽管研究进展显著，现有方法仍存在以下局限：

第一，拥塞感知与树拓扑优化耦合不足，难以实现随迭代过程动态更新的前期决策。

第二，并行化往往仅局限于单一阶段，没有对并行模板布线进行复用性拓展。



\subsection{服务布线质量提升的 VLSI 布局研究现状与进展}

\subsubsection{传统布局算法基础研究}

布局是物理设计上游核心环节，主要目标是在版图边界、密度与可制造性约束下，优化线长、拥塞、时序与功耗等多目标指标。对本文主题而言，布局研究的关键价值在于通过更优资源分布降低后续全局布线拥塞压力。经过长期发展，学术界与工业界已形成全局布局-合法化-详细布局的标准三阶段流程。

在早期研究中，模拟退火代表了随机优化路线，典型工作如 TimberWolf 系列布局器\cite{1586125}。该类方法在中小规模设计中具备较强全局搜索能力，但随设计规模增长，其计算复杂度与迭代时间快速上升，难以满足现代大规模设计需求。

为突破效率瓶颈，解析式布局成为主流技术路线。二次规划方法以 GORDIAN\cite{Kleinhans1991GORDIANVP} 和 Kraftwerk2\cite{10.1109/TCAD.2008.925783} 为代表，通过可微目标与数值优化提升收敛效率。随后，基于静电场建模与非线性优化的 EPlace\cite{10.1145/2699873} 通过泊松方程求解密度势场，在布局质量与速度上取得重要进展；RePlAce\cite{10.1109/TCAD.2018.2859220} 在密度建模与步长调节上进一步优化，显著提升了可布线性与稳健性。

在后处理阶段，合法化与详细布局用于消除重叠并减小对全局结果的扰动。Abacus\cite{10.1145/1353629.1353640} 等方法通过行内聚类与动态规划实现低位移合法化，成为主流流程中的关键组件。

\subsubsection{并行计算在布局算法中的应用演进}

随着设计规模持续扩大，串行布局算法逐渐成为设计周期瓶颈。并行计算因此成为布局加速的重要方向。早期工作主要聚焦多核 CPU 的任务并行化；近年来，GPU 异构并行成为核心趋势。DREAMPlace\cite{9122053} 将解析式布局核心算子张量化并映射到 GPU，在保持布局质量的同时显著缩短运行时间。DREAMPlace 3.0\cite{9256736} 进一步支持区域约束等复杂场景，提升了方法鲁棒性与适用范围。

与此同时，AI4EDA 路线推动了并行布局的新范式。典型工作包括基于强化学习的宏单元布局方法\cite{Mirhoseini2021AGP}，以及基于图神经网络（Graph Neural Network, GNN）的 GraphPlanner\cite{10.1145/3555804}。这类方法通常依赖 GPU/TPU 等并行平台完成训练与推理，体现了算法模型-计算架构协同优化的发展趋势。

% \subsubsection{模块化建模与布局引导优化前沿}

% 传统布局通常以标准单元为主要优化粒度，在超大规模场景下容易出现搜索空间过大、收敛缓慢等问题；同时，先布局后布线的串行流程使布局阶段难以充分感知后续布线拥塞风险。为此，面向布线质量提升的模块化建模与布局引导优化成为研究热点。

% 在可布线性引导方面，RUDY 需求估计模型\cite{10.5555/1266366.1266632} 通过快速估计布线需求辅助拥塞感知布局，为后续大量学习驱动的后端指标预测方法奠定了基础。在模块化建模方面，本文相关研究引入可微模块边界建模与多约束联合优化思想，旨在形成可直接服务于后续布局布线协同的结构化先验。

\subsubsection{布局算法现有研究局限性}

综合现有研究，布局方向仍存在以下共性问题：

第一，模块级建模能力仍有限，复杂边界与多层级约束难以统一纳入解析优化框架。

第二，解析布局算法的全流程异构适配仍不足，合法化与详细布局阶段常引入新的参数不确定性。

% 第三，布局与布线协同机制不充分，多数方法仍停留在单向预测引导，缺少闭环联合优化。

% 第四，面对更大规模设计时，内存占用、并行效率衰减与稳定收敛仍是工程落地难点。

\subsection{本章小结}

综合上述研究现状，当前布局与全局布线研究已在解析建模、并行加速与学习驱动优化方面形成丰富成果，但在模块级约束建模、细粒度增量优化方向上仍存在改进的可能。本文据此从两个方向展开：一是在初始布线及全局布线过程中引入动态拥塞感知与并行加速方法，实现质量与效率的协同提升；二是在布局阶段引入基于模块级布局规划结果的细粒度优化机制，提升布局结果的质量。本文研究覆盖2类核心环节（布线侧、布局侧）、3项主要贡献，后续实验验证将涉及31个公开布线基准与11个布局用例。
\section{本文的论文结构与章节安排}

\label{sec:arrangement}

本文共分为五章，各章节内容安排如下：

第一章为绪论。本章阐述研究背景与意义，给出本文研究内容与技术路线，综述国内外相关研究进展，并明确本文的切入点与贡献边界。

第二章为全局布线及模块级布局技术基础。本章围绕全局布线问题的基本定义、约束条件与优化目标展开，给出后续方法设计所依赖的问题建模与评价指标体系。并对模块级布局算法进行简要的介绍，为后续章节提供必要的背景信息参考。

第三章为拥塞感知初始布线工作流。本章重点介绍动态拥塞感知初始布线方法，包括候选树构建与更新机制、迭代剪枝策略以及并行求解流程，并通过实验验证其有效性。

第四章为基于模拟退火的模块级布局细粒度设计空间探索。本章在已有模块级布局规划结果基础上，利用多段模拟退火策略，研究复杂布局问题的细粒度优化方法，分析其对布局质量的影响。

第五章为总结与展望。本章总结全文工作与主要结论，讨论当前方法的局限性，并给出后续研究方向。

